\par\Needspace{6\baselineskip}\medskip
\hypertarget{native-080-article061}{}\pdfbookmark[2]{61.}{article-061}
\begin{center}61.\end{center}
Wir haben ein sehr einfaches Verfahren angegeben, mit dem sich der geozentrische Ort eines Himmelskörpers in Bezug auf die Ekliptik oder den Äquator bestimmen lässt, sowohl mit als auch ohne Berücksichtigung der Parallaxe und ebenso mit oder ohne Berücksichtigung der Nutation. Was die Nutation betrifft, so besteht der ganze Unterschied darin, ob wir die mittlere oder die wahre Lage des Äquators zugrunde legen. Im ersten Fall zählen wir die Längen daher vom mittleren Äquinoktium an, im zweiten vom wahren. Ebenso ist im ersten Fall die mittlere Schiefe der Ekliptik zu verwenden, im zweiten die wahre. Im Übrigen ist klar, dass umso mehr vorbereitende Rechenschritte nötig sind, je mehr Abkürzungen man bei der Berechnung der Koordinaten einführt. Der Vorteil des oben erläuterten Verfahrens, die Koordinaten unmittelbar aus der exzentrischen Anomalie abzuleiten, zeigt sich daher vor allem dann, wenn viele geozentrische Örter zu bestimmen sind. Ist dagegen nur ein einziger Ort oder sind nur sehr wenige Örter zu berechnen, so lohnt sich der Aufwand zur Berechnung so vieler Hilfsgrößen keineswegs. In einem solchen Fall behält man besser das übliche Verfahren bei, das aus der exzentrischen Anomalie die wahre Anomalie und den Radiusvektor ableitet. Daraus erhält man den heliozentrischen Ort in Bezug auf die Ekliptik, daraus die geozentrische Länge \SourceLeaf{081}und Breite und daraus schließlich Rektaszension und Deklination. Der Vollständigkeit halber erläutern wir auch die letzten beiden Schritte noch kurz.
